% !TEX root = ../main.tex
\chapter{Хэлэлцүүлэг}\label{ch:disc}

\section{Үр дүнгийн тайлбар}
Merge ба quick sort-д “1 thread” ба “1 thread, хувааж” хугацаа бараг тэнцүү. Учир нь хоёр хэсэгт хуваах нь
ажлыг багасгадаггүй: $2\cdot\frac{n}{2}\log\frac{n}{2} \approx n\log n$. Иймээс 1.73--1.77 хурдсалт бүхэлдээ
зэрэгцээнээс ирж байна. $n = 10^6$ үед merge sort-ын нэгтгэх үе 7.9 ms бөгөөд энэ нь ажлын
$S \approx 7.9/165.3 \approx 0.048$ хувь. Амдалын хуулиар~(\ref{eq:amdahl}) $N = 2$ үед дээд хязгаар
$1/(0.048 + 0.952/2) \approx 1.91$. Хэмжсэн 1.74 нь үүнээс бага бөгөөд зөрүүг thread үүсгэх зардал болон хоёр
цөмийн санах ойн шугамыг хуваалцдагтай холбож тайлбарлаж байна.

O($n^2$) алгоритмд хуваалт өөрөө ажлыг 2 дахин багасгадаг: $2\cdot(n/2)^2 = n^2/2$ (insertion 180 → 91 ms). Үүн
дээр зэрэгцээний ашиг нэмэгдэж, онолын нийт хурдсалт 4 болно. Insertion sort-ын 3.56 үүнд ойрхон байна. Иймээс
“O($n^2$) алгоритмд 2 thread 3.5 дахин хурдан” гэж дүгнэх нь төөрөгдөл бөгөөд ашгийн тал орчим нь алгоритмын
шинж чанараас ирж байгаа юм. Олон урсгалт хувилбар merge/quick-д $n \approx 10^5$-аас, O($n^2$)-д $n \approx 5000$-аас
эхэлж ашигтай болсон. Үүнээс бага $n$-д thread-ийн зардал давамгайлна. Жишээ нь номын жишээн дээр Sorter-L нь
Sorter-R эхлэхээс өмнө дуусчихдаг. Энэ нь алдаа биш, thread үүсгэх нь 5 элемент эрэмбэлэхээс удаан байдагтай холбоотой.

Lockstep горимд хоёр sorter L, R, L, R гэж хатуу ээлжилсэн. Эхэндээ үүнийг OS-ийн scheduler шийддэг гэж үзсэн боловч
шалтгаан нь fair semaphore-ийн FIFO дараалал байв. Харин эхний permit-ийг аль sorter авахыг scheduler шийддэг.

\section{Хязгаарлалт}\label{sec:limitations}
Хэмжилтийг 2 цөмтэй виртуал машин дээр хийсэн тул JIT ба GC-ийн туслах thread-үүд sorter-уудтай цөмийн төлөө
өрсөлддөг. Bubble sort-ын 3 thread-ийн хугацаа JVM-ээс хамаарч 67--145 ms хооронд хэлбэлзсэн бөгөөд түүний
зэрэгцээний хурдсалт 2.68 нь 2 цөмийн физик хязгаараас давсан. Иймд энэ утгыг JIT-ийн нөлөө гэж үзэж итгэл муутай
гэж тооцсон. 1 ms-ээс бага хугацаа шуугиантай. Windows-ийн \code{.bat} скриптийг бодит Windows машин дээр туршаагүй.

\section{Сайжруулах боломж}
$k$ sorting thread бүхий мод хэлбэрийн нэгтгэл нь олон цөмийг бүрэн ашиглана. Нэгтгэх үеийг хоёр thread-д хуваавал
(нэг нь эхнээс, нөгөө нь төгсгөлөөс) Амдалын $S$ багасна. Thread pool болон босго утга (жижиг $n$-д дан thread)
ашиглавал жижиг өгөгдлийн зардал буурна. \code{Arrays.parallelSort} ч мөн ийм босго ашигладаг~\cite{javadoc-arrays}.

\section{GitHub ба AI ашигласан тухай}
\code{22dbddf} ба \code{a8916d6} commit-ууд ижил мессежтэй. Учир нь web UI-г шалгалт дуусахаас өмнө commit хийсэн
тул засвар нь дараагийн commit-д орсон. Хуваалцсан түүхийг \code{force push}-оор өөрчлөхгүйн тулд хэвээр үлдээв.

Төслийн код, тест, гарын авлага, тайлангийн нооргийг бэлтгэхэд Claude хиймэл оюун ухаан ашигласан.
Бүх \code{git} үйлдлийг багийн гишүүд өөрсдөө хийсэн. Гишүүд AI-ийн үүсгэсэн материалыг уншиж, засварлаж, програмыг
өөрсдийн компьютер дээр ажиллуулж шалгасан. Тайлан дахь тоо бүр програмыг бодитоор ажиллуулж гарсан утга юм.
